% part: intuitionistic-logic
% chapter: soundness-completeness
% section: canonical-model

\documentclass[../../../include/open-logic-section]{subfiles}

\begin{document}

\olfileid{int}{sc}{mod}

\olsection{Model Kanonik}

Donya ing model kita bakal awujud urutan winates~$\sigma$ saka
wilangan asli, yaiku $\sigma \in \Nat^*$. $\Nat^*$ ditetepake sacara
induktif dening:
\begin{enumerate}
\item $\emptyseq \in \Nat^*$.
\item Yen $\sigma \in \Nat^*$ lan $n \in \Nat$, mula $\sigma.n \in
  \Nat^*$ (kanthi $\sigma.n$ yaiku $\sigma \concat \tuple{n}$ lan
  $\sigma \concat \sigma'$ yaiku konkatenasi saka $\sigma$ lan $\sigma'$).
\item Ora ana liyane kang ana ing $\Nat^*$.
\end{enumerate}
Mula kita bisa nggunakake $\Nat^*$ kanggo menehi definisi induktif.

Upamane $\tuple{!B_0, !C_0}$, $\tuple{!B_1, !C_1}$, \dots, iku
enumerasi kabeh pasangan !!{formula}s. Yen diwenehi himpunan
!!{formula}s~$\Delta$, tetepa $\Delta(\sigma)$ nganggo induksi
kaya mangkene:
\begin{enumerate}
\item $\Delta(\emptyseq) = \Delta$
\item $\Delta(\sigma.n) = {}$
  \[
  \begin{cases}
    (\Delta(\sigma) \cup \{!B_n\})^* &
    \text{yen $\Delta(\sigma) \cup \{!B_n\} \Proves/ !C_n$} \\
    \Delta(\sigma) & \text{yen ora}
  \end{cases}
  \]
\end{enumerate}
Ing kene, $(\Delta(\sigma) \cup \{!B_n\})^*$ tegese himpunan prim saka
!!{formula}s kang ana miturut \olref[lin]{lem:lindenbaum}, kanthi
lema iku ditrapake marang himpunan $\Delta(\sigma) \cup \{!B_n\}$
lan !!{formula}~$!C_n$. Miturut definisi iki, yen
$\Delta(\sigma) \cup \{!B_n\} \Proves/ !C_n$, mula
$\Delta(\sigma.n) \Proves !B_n$ lan $\Delta(\sigma.n) \Proves/ !C_n$.
Uga $\Delta(\sigma) \subseteq \Delta(\sigma.n)$ kanggo sembarang~$n$.
Yen $\Delta$ prim, mula $\Delta(\sigma)$ prim kanggo kabeh~$\sigma$.

\begin{defn}\ollabel{defn:canonical-model}
  Upamane $\Delta$ prim. Mula \emph{model kanonik} $\mModel{M(\Delta)}$
  kanggo $\Delta$ ditetepake dening:
  \begin{enumerate}
  \item $W = \Nat^*$, himpunan urutan winates saka wilangan asli.
  \item $R$ iku urutan parsial kanthi $R\sigma\sigma'$ yen lan mung yen
    $\sigma$ iku segmen wiwitan saka~$\sigma'$ (yaiku
    $\sigma' = \sigma \concat \sigma''$ kanggo sawenehing urutan~$\sigma''$).
  \item $V(p) = \Setabs{\sigma}{p \in \Delta(\sigma)}$.
  \end{enumerate}
\end{defn}

Gampang dipriksa yen $R$ pancen urutan parsial. Syarat monotonisitas
tumrap~$V$ uga kasembadan. Amarga
$\Delta(\sigma) \subseteq \Delta(\sigma.n)$, kita entuk
$\Delta(\sigma) \subseteq \Delta(\sigma')$ saben $R\sigma\sigma'$
nganggo induksi tumrap dawa urutan kang ditambahake marang~$\sigma$.

\end{document}
